\documentclass[11pt,a4paper]{article}

% --- Packages ---
\usepackage[utf8]{inputenc}
\usepackage[T1]{fontenc}
\usepackage{amsmath,amssymb,amsthm}
\usepackage{physics}
\usepackage{geometry}
\usepackage{enumitem}
\usepackage{hyperref}
\usepackage{fancyhdr}
\usepackage{mathtools}

\geometry{margin=1in}
\pagestyle{fancy}
\fancyhf{}
\rhead{Quantum Spin Lecture Notes}
\lhead{Lecture 1: Introduction to Quantum Spin}
\cfoot{\thepage}

% --- Theorem Environments ---
\theoremstyle{definition}
\newtheorem{definition}{Definition}[section]
\newtheorem{example}{Example}[section]
\theoremstyle{plain}
\newtheorem{theorem}{Theorem}[section]
\newtheorem{proposition}{Proposition}[section]
\theoremstyle{remark}
\newtheorem{remark}{Remark}[section]

\title{\textbf{Quantum Spin (1): Introduction}}
\author{Lecture Notes on Quantum Spin}
\date{}

\begin{document}
\maketitle
\tableofcontents
\newpage

% ============================================================
\section{Motivation}
% ============================================================

Quantum spin for spin-$\frac{1}{2}$ particles (electrons, protons, neutrons) is often described through analogies: ``spin is \emph{like} the particle is actually spinning, but not really,'' or ``spin is \emph{like} an angular momentum vector, but not really.'' These analogies, while intuitive, obscure the mathematical simplicity of quantum spin.

The goal of these notes is to describe \emph{what spin is} --- not merely what it is like. The only prerequisite is familiarity with complex numbers.

% ============================================================
\section{Review of Complex Numbers}
% ============================================================

\begin{definition}[Complex Numbers]
    The set of \textbf{complex numbers} $\mathbb{C}$ consists of all numbers of the form $z = a + bi$, where $a, b \in \mathbb{R}$ and $i^2 = -1$. Here $a = \operatorname{Re}(z)$ is the \textbf{real part} and $b = \operatorname{Im}(z)$ is the \textbf{imaginary part}.
\end{definition}

\begin{definition}[Complex Conjugate]
    If $z = a + bi$, the \textbf{complex conjugate} is $\bar{z} = a - bi$.
\end{definition}

\begin{definition}[Magnitude Squared]
    The \textbf{magnitude squared} of $z$ is
    \[
        |z|^2 \coloneqq \bar{z}\, z = (a - bi)(a + bi) = a^2 + b^2.
    \]
    Since $a, b \in \mathbb{R}$, we have $|z|^2 \geq 0$, with equality if and only if $z = 0$.
\end{definition}

% ============================================================
\section{Vectors in $\mathbb{R}^3$: A Quick Review}
% ============================================================

\subsection{Components and Basis Vectors}

A familiar vector $\vec{L} \in \mathbb{R}^3$ is specified by three real components:
\[
    \vec{L} = \begin{pmatrix} L_x \\ L_y \\ L_z \end{pmatrix} = L_x\,\hat{\imath} + L_y\,\hat{\jmath} + L_z\,\hat{k},
\]
where the standard basis vectors are
\[
    \hat{\imath} = \begin{pmatrix}1\\0\\0\end{pmatrix}, \quad
    \hat{\jmath} = \begin{pmatrix}0\\1\\0\end{pmatrix}, \quad
    \hat{k} = \begin{pmatrix}0\\0\\1\end{pmatrix}.
\]

\subsection{The Dot Product}

The \textbf{dot product} of $\vec{v}, \vec{u} \in \mathbb{R}^3$ is
\[
    \vec{v} \cdot \vec{u} = v_x u_x + v_y u_y + v_z u_z \in \mathbb{R}.
\]
Key properties:
\begin{itemize}
    \item Symmetry: $\vec{v} \cdot \vec{u} = \vec{u} \cdot \vec{v}$.
    \item Basis orthonormality: $\hat{\imath} \cdot \hat{\imath} = 1$, $\hat{\imath} \cdot \hat{\jmath} = 0$, etc.
    \item Component extraction: $v_x = \vec{v} \cdot \hat{\imath}$, $v_y = \vec{v} \cdot \hat{\jmath}$, $v_z = \vec{v} \cdot \hat{k}$.
    \item Geometric meaning: $\vec{v} \cdot \vec{u} = |\vec{v}|\,|\vec{u}|\cos\theta$, measuring how much two vectors ``overlap.''
\end{itemize}

% ============================================================
\section{Quantum Spin Vectors}
% ============================================================

\subsection{Two Key Differences from Classical Vectors}

Quantum spin vectors differ from ordinary $\mathbb{R}^3$ vectors in two fundamental ways:
\begin{enumerate}[label=\textbf{\arabic*.}]
    \item \textbf{Dimension:} Spin vectors are \textbf{2-dimensional}, not 3-dimensional.
    \item \textbf{Components:} Spin vectors have \textbf{complex} components, not real.
\end{enumerate}

That is, quantum spin states are elements of $\mathbb{C}^2$, not $\mathbb{R}^3$.

\subsection{Notation: Kets and Basis States}

In Dirac notation, the basis vectors for spin-$\frac{1}{2}$ are written as \textbf{kets}:
\[
    \ket{\uparrow} = \begin{pmatrix} 1 \\ 0 \end{pmatrix}, \qquad
    \ket{\downarrow} = \begin{pmatrix} 0 \\ 1 \end{pmatrix},
\]
read as ``spin up'' and ``spin down.'' A general spin state is
\[
    \ket{\psi} = \alpha\ket{\uparrow} + \beta\ket{\downarrow} = \begin{pmatrix} \alpha \\ \beta \end{pmatrix} \in \mathbb{C}^2,
\]
where $\alpha, \beta \in \mathbb{C}$.

\begin{remark}
    Despite the unfamiliar notation, the mathematics is essentially the same as for ordinary vectors --- just with two complex components instead of three real ones.
\end{remark}

% ============================================================
\section{The Inner Product on $\mathbb{C}^2$}
% ============================================================

\subsection{Definition}

The analog of the dot product for spin vectors involves complex conjugation:

\begin{definition}[Inner Product on $\mathbb{C}^2$]
    For $\ket{\psi_1} = \begin{pmatrix} \alpha_1 \\ \beta_1 \end{pmatrix}$ and $\ket{\psi_2} = \begin{pmatrix} \alpha_2 \\ \beta_2 \end{pmatrix}$, the \textbf{inner product} is
    \[
        \braket{\psi_2}{\psi_1} = \bar{\alpha}_2\,\alpha_1 + \bar{\beta}_2\,\beta_1 \in \mathbb{C}.
    \]
\end{definition}

\subsection{Properties}

\begin{itemize}
    \item The inner product is a single complex number (not a vector).
    \item \textbf{Conjugate symmetry} (not full symmetry):
    \[
        \braket{\psi_1}{\psi_2} = \overline{\braket{\psi_2}{\psi_1}}.
    \]
    \item The inner product of a state with itself is real and non-negative:
    \[
        \braket{\psi}{\psi} = |\alpha|^2 + |\beta|^2 \geq 0.
    \]
\end{itemize}

% ============================================================
\section{Dual Vectors (Bras)}
% ============================================================

\begin{definition}[Dual Vector / Bra]
    Given a ket $\ket{\psi} = \begin{pmatrix} \alpha \\ \beta \end{pmatrix}$, the \textbf{dual vector} (or \textbf{bra}) is obtained by complex conjugating the components and transposing to a row vector:
    \[
        \bra{\psi} = \begin{pmatrix} \bar{\alpha} & \bar{\beta} \end{pmatrix}.
    \]
\end{definition}

The inner product can then be written as a matrix product:
\[
    \braket{\psi_2}{\psi_1} = \begin{pmatrix} \bar{\alpha}_2 & \bar{\beta}_2 \end{pmatrix} \begin{pmatrix} \alpha_1 \\ \beta_1 \end{pmatrix} = \bar{\alpha}_2 \alpha_1 + \bar{\beta}_2 \beta_1.
\]

\begin{remark}[Origin of the Names]
    This notation was invented by Paul Dirac. A ket $\ket{\psi}$ is a vector; a bra $\bra{\psi}$ is a dual vector. Placing them together, $\braket{\psi_2}{\psi_1}$, forms a \textbf{bra-ket} (bracket) --- the inner product.
\end{remark}

In matrix language:
\begin{itemize}
    \item A ket is a $2 \times 1$ matrix (column vector).
    \item A bra is a $1 \times 2$ matrix (row vector).
    \item Their product $\braket{\cdot}{\cdot}$ is a $1 \times 1$ matrix (a number).
\end{itemize}

% ============================================================
\section{Measurements and Probability}
% ============================================================

\subsection{The Experimental Setup}

An experimenter can ask one question about an electron's spin: \emph{Is the spin up or down with respect to a particular axis?}

\begin{itemize}
    \item There are only \textbf{two possible outcomes}: spin up or spin down.
    \item A theorist can only predict the \textbf{probabilities} of each outcome --- not which one will occur with certainty.
\end{itemize}

\subsection{Probability Rules}

Given a state $\ket{\psi} = \begin{pmatrix} \alpha \\ \beta \end{pmatrix}$, the measurement probabilities along the $z$-axis are:
\begin{align}
    P(\text{spin up along } z) &= |\alpha|^2 = \bigl|\braket{\uparrow_z}{\psi}\bigr|^2, \\[4pt]
    P(\text{spin down along } z) &= |\beta|^2 = \bigl|\braket{\downarrow_z}{\psi}\bigr|^2.
\end{align}

Special cases:
\begin{itemize}
    \item $\ket{\psi} = \ket{\uparrow_z} = \begin{pmatrix}1\\0\end{pmatrix}$: always measured spin up ($100\%$).
    \item $\ket{\psi} = \ket{\downarrow_z} = \begin{pmatrix}0\\1\end{pmatrix}$: always measured spin down ($100\%$).
\end{itemize}

\subsection{Normalization}

Since probabilities must sum to 1:
\[
    |\alpha|^2 + |\beta|^2 = 1 \quad \Longleftrightarrow \quad \braket{\psi}{\psi} = 1.
\]
A state satisfying this condition is called \textbf{properly normalized} or a \textbf{physical state}.

% ============================================================
\section{Spin States Along Other Axes}
% ============================================================

The states that are always spin up or down along each axis are:

\begin{center}
\renewcommand{\arraystretch}{1.6}
\begin{tabular}{c|c|c}
    \textbf{Axis} & $\ket{\uparrow}$ & $\ket{\downarrow}$ \\ \hline
    $x$ & $\dfrac{1}{\sqrt{2}}\begin{pmatrix}1\\1\end{pmatrix}$ & $\dfrac{1}{\sqrt{2}}\begin{pmatrix}1\\-1\end{pmatrix}$ \\[8pt]
    $y$ & $\dfrac{1}{\sqrt{2}}\begin{pmatrix}1\\i\end{pmatrix}$ & $\dfrac{1}{\sqrt{2}}\begin{pmatrix}1\\-i\end{pmatrix}$ \\[8pt]
    $z$ & $\begin{pmatrix}1\\0\end{pmatrix}$ & $\begin{pmatrix}0\\1\end{pmatrix}$
\end{tabular}
\end{center}

The factors of $1/\sqrt{2}$ ensure normalization: $\braket{\psi}{\psi} = 1$.

\subsection{General Measurement Rule}

For \emph{any} axis $a \in \{x, y, z\}$:
\[
    P(\text{spin up along } a) = \bigl|\braket{\uparrow_a}{\psi}\bigr|^2, \qquad
    P(\text{spin down along } a) = \bigl|\braket{\downarrow_a}{\psi}\bigr|^2.
\]

\subsection{Orthogonality}

Along every axis, the spin-up and spin-down states are orthogonal:
\[
    \braket{\uparrow_a}{\downarrow_a} = 0 \quad \text{for } a \in \{x, y, z\}.
\]

\subsection{Context-Dependent Definite Spin}

Whether a particle has a definite spin depends on which axis you measure. For example, $\ket{\uparrow_x}$ has definite spin along $x$ (always up), but along $z$ it yields:
\[
    P(\text{up along } z) = \left|\frac{1}{\sqrt{2}}\right|^2 = \frac{1}{2} = 50\%.
\]
We say $\ket{\uparrow_x}$ is in a \textbf{superposition} of the $z$-axis eigenstates:
\[
    \ket{\uparrow_x} = \frac{1}{\sqrt{2}}\ket{\uparrow_z} + \frac{1}{\sqrt{2}}\ket{\downarrow_z}.
\]

% ============================================================
\section{Worked Example}
% ============================================================

Consider the state
\[
    \ket{\psi} = \frac{3}{5}\ket{\uparrow_z} + \frac{4i}{5}\ket{\downarrow_z} = \begin{pmatrix} 3/5 \\ 4i/5 \end{pmatrix}.
\]

\textbf{Normalization check:}
\[
    \braket{\psi}{\psi} = \left(\frac{3}{5}\right)^2 + \left(\frac{4}{5}\right)^2 = \frac{9}{25} + \frac{16}{25} = 1. \quad \checkmark
\]

\textbf{Probabilities along the $z$-axis:}
\[
    P(\uparrow_z) = \left|\frac{3}{5}\right|^2 = \frac{9}{25} = 36\%, \qquad
    P(\downarrow_z) = \left|\frac{4}{5}\right|^2 = \frac{16}{25} = 64\%.
\]

\textbf{Probabilities along the $x$-axis:}
\begin{align}
    P(\uparrow_x) &= \left|\braket{\uparrow_x}{\psi}\right|^2 = \left|\frac{1}{\sqrt{2}}\left(\frac{3}{5} + \frac{4i}{5}\right)\right|^2 = \frac{1}{2}\left(\frac{9}{25} + \frac{16}{25}\right) = 50\%, \\
    P(\downarrow_x) &= 50\%.
\end{align}

\textbf{Probabilities along the $y$-axis:}
\begin{align}
    P(\uparrow_y) &= \left|\braket{\uparrow_y}{\psi}\right|^2 = \left|\frac{1}{\sqrt{2}}\left(\frac{3}{5} + \frac{4i}{5}\cdot(-i)\right)\right|^2 = \left|\frac{1}{\sqrt{2}}\cdot\frac{7}{5}\right|^2 = \frac{49}{50} = 98\%, \\
    P(\downarrow_y) &= \frac{1}{50} = 2\%.
\end{align}

This state is nearly aligned with the $+y$ direction.

% ============================================================
\section{Quantum Measurement and Collapse}
% ============================================================

\subsection{The Collapse Postulate}

After a measurement is performed, the state \textbf{collapses} onto the measured outcome:

\begin{itemize}
    \item If the initial state is $\ket{\psi} = \alpha\ket{\uparrow} + \beta\ket{\downarrow}$ and spin up is measured, then the post-measurement state becomes $\ket{\uparrow}$.
    \item Subsequent identical measurements will yield spin up with $100\%$ probability.
\end{itemize}

\subsection{Collapse Prevents State Determination}

Collapse prevents an experimenter from determining the quantum state $(\alpha, \beta)$ of a single electron with arbitrary precision. Each measurement alters the state, making it impossible to reconstruct the original state from repeated measurements on the \emph{same} particle.

\subsection{Intermediate Measurements Disturb the State}

Consider an electron initially in state $\ket{\uparrow_z}$:
\begin{enumerate}
    \item Measuring along $z$: $100\%$ spin up.
    \item Now insert an intermediate $x$-axis measurement (yielding up or down, each $50\%$), \emph{then} measure along $z$ again.
    \item Result: $50\%$ spin up, $50\%$ spin down along $z$ --- the intermediate measurement has destroyed the original information.
\end{enumerate}

% ============================================================
\section{Global Phase Invariance}
% ============================================================

If $|z| = 1$ (i.e., $z = e^{i\theta}$ for some $\theta \in \mathbb{R}$), then the states $\ket{\psi}$ and $z\ket{\psi}$ are \textbf{physically indistinguishable}: all measurement probabilities are identical.

\begin{proposition}
    For any $z \in \mathbb{C}$ with $|z| = 1$ and any axis $a$:
    \[
        \bigl|\braket{\uparrow_a}{z\psi}\bigr|^2 = |z|^2 \bigl|\braket{\uparrow_a}{\psi}\bigr|^2 = \bigl|\braket{\uparrow_a}{\psi}\bigr|^2.
    \]
\end{proposition}

The factor $z = e^{i\theta}$ is called an \textbf{overall (global) phase} and is unobservable.

% ============================================================
\section{Spin Quantum Number and Angular Momentum}
% ============================================================

The spin quantum number $s$ of a particle determines the dimension of its state space:
\[
    \dim = 2s + 1.
\]

\begin{center}
\renewcommand{\arraystretch}{1.3}
\begin{tabular}{c|c|c|c}
    $s$ & $2s+1$ & State space & Angular momentum values ($m_s \hbar$) \\ \hline
    $0$ & $1$ & $\mathbb{C}^1$ & $0$ \\
    $1/2$ & $2$ & $\mathbb{C}^2$ & $-\hbar/2,\; +\hbar/2$ \\
    $1$ & $3$ & $\mathbb{C}^3$ & $-\hbar,\; 0,\; +\hbar$ \\
    $3/2$ & $4$ & $\mathbb{C}^4$ & $-3\hbar/2,\; -\hbar/2,\; +\hbar/2,\; +3\hbar/2$
\end{tabular}
\end{center}

Angular momentum is quantized in units of $\hbar$: consecutive values always differ by exactly $\hbar$.

\end{document}
